\chapter{Conclusions and Future Work}\label{ch:conclusion}

This thesis has developed \textit{AntCalc}, a computational framework for the construction and
analytic integration of QCD antenna functions. The principal objective was to automate a technically
demanding component of antenna subtraction: the derivation of the antenna functions from the
corresponding matrix elements, their reduction to suitable integral representations, and their
analytic integration over the unresolved phase space. \textit{AntCalc} was developed from the ground
up in this work and provides a common framework in which these different stages can be carried out
consistently.

The main result is the implementation of the complete set of massless final--final quark--antiquark
antenna functions considered in this work through NNLO. This includes the tree-level, one-loop and
two-loop two-parton contributions, the three-parton tree-level and real--virtual contributions, and
the four-parton double-real antennae with both gluonic and fermionic final states. The antenna
functions are available both in their unintegrated form, required for the construction of local
subtraction terms, and in analytically integrated form, required for the cancellation of the explicit
infrared poles of virtual contributions. Their integration required the reduction of the
corresponding phase-space and loop integrals to master integrals, using IBP reduction where
appropriate, followed by the substitution of their analytic expressions and expansion in the
dimensional-regularisation parameter $\epsilon$.

The results were subjected to several complementary validation procedures. The unintegrated antenna
functions were compared with established results in the literature, and their infrared behaviour was
tested in the appropriate unresolved limits. In particular, the massless $A_3^0$ antenna was used to
demonstrate explicitly the recovery of the soft eikonal factor and the relevant Altarelli--Parisi
splitting functions. At the integrated level, the complete massless antenna set was subjected to the
global validation presented in Chapter~\ref{ch:validation}. When assembled into the massless hadronic
$R$-ratio, the infrared poles of the separate real and virtual contributions cancel at NLO and NNLO,
and the remaining finite coefficients reproduce the independently known result. This provides a
particularly stringent test of the relative normalisation, colour decomposition, and integrated
infrared structure of the antenna set.

This validation is important in view of the intended application of these functions. Antenna
functions are not isolated observables, but universal building blocks used to construct subtraction
terms for higher-order QCD calculations. A consistent NNLO subtraction calculation requires the
appropriate unintegrated antennae to reproduce the local unresolved behaviour of real-radiation
matrix elements, while their integrated counterparts must provide the infrared poles required to
cancel those of the corresponding virtual contributions. The local and global checks performed in
this thesis therefore test precisely the properties required for the antenna functions to be used
within the antenna-subtraction framework.

Beyond the complete massless quark--antiquark set, this work has also explored the extension of
\textit{AntCalc} to massive antenna functions through the implementation of the massive $A_3^0$
antenna. This required a separate massive integral family and a transformation between the
master-integral basis produced by the IBP reduction and that used for the known analytic results.
Although this route has not yet undergone the same level of global validation as the massless antenna
set, it demonstrates that the architecture of \textit{AntCalc} is not intrinsically restricted to
massless kinematics and provides a starting point for a systematic extension to massive antenna
functions.

\section*{Future Work}

The results presented in this thesis provide several natural directions for the further development of
\textit{AntCalc}. From the physics perspective, the most immediate extension is the completion of the
massive quark--antiquark final--final antenna set. The massive $A_3^0$ antenna implemented in this
work provides a first demonstration that the framework developed for the massless case can be extended
to massive kinematics. Completing this sector would require the corresponding massive four-parton
antennae, $A_4^0$, $B_4^0$, and $C_4^0$, followed by the massive real--virtual and double-virtual
antennae $A_3^1$ and $A_2^2$. Importantly, this extension would go beyond reproducing antenna
functions that are already analytically known: some of the required massive antenna functions remain
to be calculated. The framework developed and validated in this thesis therefore provides a possible
starting point for obtaining genuinely new ingredients for massive NNLO antenna subtraction.

A broader extension would be to include antenna functions associated with other pairs of hard
radiators. The present implementation focuses on quark--antiquark antennae, while a general
antenna-subtraction calculation also requires quark--gluon and gluon--gluon antenna functions,
conventionally organised into the $D$- and $E$-type and $F$-, $G$-, and $H$-type families,
respectively. Extending \textit{AntCalc} to these antenna families would broaden the range of partonic
processes for which the required subtraction ingredients could be constructed within a common
framework. Similarly, extending the present final--final implementation to initial--final and
initial--initial kinematics would be necessary for applications involving coloured partons in the
initial state. In the longer term, the same framework could provide a basis for investigating the
antenna functions required at perturbative orders beyond NNLO, in particular at N$^3$LO.

These extensions would substantially increase the algebraic complexity and computational cost of both
the construction and integration stages. They therefore also motivate the further development of the
computational architecture of \textit{AntCalc}. In its current form, the package is implemented
entirely within \textit{Wolfram Mathematica}. This environment has provided a flexible setting in which
to develop, test, and validate the symbolic workflow presented in this thesis, but the more involved
calculations already reveal limitations in execution time and memory usage. For example, the complete
\texttt{BuildRRatio} calculation currently requires approximately one and a half to two hours, with
its execution becoming sensitive to the available memory. While this does not affect the results
obtained in the present work, it indicates that a more scalable implementation would be desirable
before substantially larger antenna families and integral reductions are attempted.

One possible architecture for this development is illustrated in Fig.~\ref{fig:next}. Rather than
replacing the high-level workflow established in this thesis, the objective would be to make its
implementation more modular and to delegate computationally intensive stages to specialised tools. A
\textit{Python}-based layer could coordinate the workflow, with the process-dependent profiles stored
in a portable, machine-readable format such as \texttt{JSON}. \textit{FeynArts} and \textit{FeynCalc}
could continue to provide diagram generation and amplitude manipulation, respectively, while large
symbolic manipulations could be delegated to \textit{FORM}. Similarly, the IBP-reduction stage could
be transferred from \textit{LiteRed2} to \textit{Kira}, while \textit{Mathematica} could be retained
for symbolic operations for which it remains advantageous. The purpose of this restructuring would
therefore not simply be to improve the performance of calculations that \textit{AntCalc} can already
perform, but to make the framework sufficiently scalable for the more demanding physics applications
outlined above.

\begin{figure}[h!]
\centering
\begin{tikzpicture}[
  font=\small,
  >=Stealth,
  stage/.style={
    draw=black!65,
    rounded corners=2pt,
    fill=white,
    align=center,
    minimum height=10mm,
    text width=19mm,
    inner sep=3pt
  },
  profile/.style={
    stage,
    fill=violet!7,
    text width=34mm,
    minimum height=12mm
  },
  public/.style={
    stage,
    fill=blue!7,
    text width=24mm
  },
  flow/.style={->, line width=.75pt},
  config/.style={->, dashed, gray!75, line width=.6pt}
]

\tikzset{
  % Basic middle arrow
  mid arrow/.style={
    decoration={
      markings,
      mark=at position 0.5 with {\arrow{Triangle}}
    },
    postaction={decorate}
  },
  % Flexible version with custom position & arrow tip
  mid arrow custom/.style 2 args={
    decoration={
      markings,
      mark=at position #1 with {\arrow{#2}}
    },
    postaction={decorate}
  }
}

% 1st card
\filldraw[black] (2.5,5) circle (2pt) node[align=center, anchor=south, yshift = 2pt]{Diagram \\ \& Amplitude};
\filldraw[black] (2.5,3) circle (2pt) node[anchor=north, yshift = -2pt]{Algebra};
\filldraw[black] (5.5,5) circle (2pt) node[align=center, anchor=south, yshift = 2pt]{Laurent \\ Expansion};
\filldraw[black] (5.5,3) circle (2pt) node[align=center, anchor=north, yshift = -2pt]{IBP \\ Reduction};

\draw[mid arrow, thick] (2.5,5) -- (2.5,3);
\draw[mid arrow, thick] (2.5,3) -- (5.5,3);
\draw[mid arrow, thick] (5.5,3) -- (5.5,5);

\draw[dashed, black!40] (\textwidth/4-0.125cm,7) -- (2.5,3);
\draw[dashed, black!40] (\textwidth/4-0.125cm,7) -- (2.5,5);
\draw[dashed, black!40] (\textwidth/4-0.125cm,7) -- (5.5,3);
\draw[dashed, black!40] (\textwidth/4-0.125cm,7) -- (5.5,5);

\filldraw[black] (\textwidth/4-0.125cm,7) circle (2pt) node[anchor=south, yshift = 2pt]{Orchestrator};

% 2nd card
\filldraw[black] (10.125,4.5) circle (2pt) node[align=center, anchor=south, yshift = 2pt]{Diagram \\ \& Amplitude};
\filldraw[black] (10.4375,1.8) circle (2pt) node[align=center, anchor=north, yshift = -2pt]{Algebra};
\filldraw[black] (13.8125,1.8) circle (2pt) node[align=center, anchor=north, yshift = -2pt]{IBP \\ Reduction};
\filldraw[black] (14.125,4.5) circle (2pt) node[align=center, anchor=south, yshift = 2pt]{Laurent \\ Expansion};

\draw[mid arrow, thick] (10.125,4.5) -- (10.4375,1.8);
\draw[mid arrow, thick] (10.4375,1.8) -- (13.8125,1.8);
\draw[mid arrow, thick] (13.8125,1.8) -- (14.125,4.5);

\draw[dashed, black!40] (12.125cm,7.5) -- (10.125,4.5);
\draw[dashed, black!40] (12.125cm,7.5) -- (10.4375,1.8);
\draw[dashed, black!40] (12.125cm,7.5) -- (13.8125,1.8);
\draw[dashed, black!40] (12.125cm,7.5) -- (14.125,4.5);

\begin{pgfonlayer}{background}
  % Draws a light blue rectangle behind a specific region
  \fill[black!10, rounded corners] (0,0) rectangle (7.75,9);
  \fill[red!10, rounded corners] (1,1) rectangle (6.75,8);

  \fill[black!10, rounded corners] (8.25,0) rectangle (16,9);
  \fill[blue!15, rounded corners] (10.75,6.5) rectangle (13.5,8.5);
  \fill[red!10, rounded corners] (8.75,3.5) rectangle (11.5,6);
  \fill[orange!15, rounded corners] (9,2.8) rectangle (11.875,0.3);
  \fill[cyan!15, rounded corners] (12.375,2.8) rectangle (15.25,0.3);
  \fill[red!10, rounded corners] (12.75,3.5) rectangle (15.5,6);
\end{pgfonlayer}

% labels
\node[label={[text=gray]right:\textit{Current}}] at (0,8.5) {};
\node[label={[text=red!40]right:\textit{Mathematica}}] at (1,1.3) {};

\node[label={[text=gray]right:\textit{Planned}}] at (8.5,8.5) {};
\node[label={[text=blue!40]right:\textit{Python}}] at (10.75,6.8) {};
\node[label={[text=red!40]right:\textit{Mathematica}}] at (8.75,3.8) {};
\node[label={[text=orange!60]right:\textit{FORM}}] at (9,0.6) {};
\node[label={[text=cyan!60]right:\textit{C++}}] at (12.375,0.6) {};
\node[label={[text=red!40]right:\textit{Mathematica}}] at (12.75,3.8) {};

\filldraw[black] (12.125cm,7.5) circle (2pt) node[anchor=south, yshift = 2pt]{Orchestrator};

\end{tikzpicture}
\caption{\centering Diagram showing the current and planned workflow steps relative to their programming languages.}
\label{fig:higher-level}
\end{figure}


The present work thus provides both a validated set of massless quark--antiquark antenna functions
through NNLO and a framework from which the calculation of further antenna functions can be pursued.
The extension to the massive $A_3^0$ antenna demonstrates the first step beyond the massless sector,
while the remaining massive antennae provide a particularly well-defined direction in which the
methods developed here can be applied to calculations that are not yet complete in the literature.
Extending the framework to additional radiator configurations and kinematics would further broaden its
applicability to antenna subtraction. In this sense, the developments proposed above are not separate
from the work presented in this thesis, but represent a continuation of its central objective: to make
the construction and analytic integration of the antenna functions required for higher-order QCD
calculations more systematic, reproducible, and computationally accessible.
